% © 2026 Ing. David Jaroš — CC BY-NC-ND 4.0
%
% This work is licensed under a Creative Commons Attribution-NonCommercial-NoDerivatives
% 4.0 International License (CC BY-NC-ND 4.0).
%
% License History: Earlier drafts (up to v0.3) were released under CC BY 4.0.
% From v0.4 onward, all material is released under CC BY-NC-ND 4.0 to protect
% the integrity of the theoretical work during ongoing academic development.
%
% See LICENSE.md for full license text.

% Mirror Sector: Two Independent Vacua in UBT
% Status: [MOTIVATED CONJECTURE — two independent algebraic sectors]
% Date: 2026-03-07

\ifdefined\INCLUDEMODE
\else
  \documentclass[12pt]{article}
  \usepackage[a4paper, margin=2.5cm]{geometry}
  \usepackage{amsmath, amssymb, amsthm}
  \usepackage{hyperref}
  \usepackage{xcolor}
  \usepackage{mdframed}
  \usepackage{booktabs}
  \usepackage{tcolorbox}

  \newtheorem{theorem}{Theorem}[section]
  \newtheorem{lemma}[theorem]{Lemma}
  \newtheorem{proposition}[theorem]{Proposition}
  \theoremstyle{definition}
  \newtheorem{definition}[theorem]{Definition}
  \theoremstyle{remark}
  \newtheorem{remark}[theorem]{Remark}

  \newmdenv[backgroundcolor=yellow!20, linecolor=orange, linewidth=1.5pt]{gapbox}
  \newmdenv[backgroundcolor=green!15, linecolor=green!60!black, linewidth=1.5pt]{resultbox}
  \newmdenv[backgroundcolor=blue!10, linecolor=blue!60, linewidth=1.5pt]{conjecturebox}

  \tcbuselibrary{skins,breakable}
  \newtcolorbox{statusbox}[1]{
    colback=yellow!8!white, colframe=orange!60!black,
    title={\textbf{#1}}, fonttitle=\bfseries, breakable
  }

  \title{\textbf{Mirror Sector: Two Independent Vacua in UBT}\\
         \large $n^* = 137$ and $n^* = 139$ as Parallel Solutions}
  \author{David Jaroš}
  \date{2026-03-07}

  \begin{document}
  \maketitle

  \begin{abstract}
  The UBT effective potential $V_{\rm eff}(n) = A\cdot n^2 - B\cdot n\ln(n)$
  has a unique global minimum at $n^* = 137$ when calibrated to our sector.
  We show that the twin prime $n^* = 139$ is \emph{not} a local minimum
  of the same potential, but rather the global minimum of a parallel branch
  of the theory with a slightly different loop-correction coefficient
  $B' = B_{139} = 2A\cdot 139/(\ln 139 + 1)$, differing from $B_{137}$
  by only $\Delta B/B = 1.21\%$.

  The mirror sector with $n^* = 139$ is therefore a \textbf{fully autonomous
  vacuum} --- not a metastable excitation.  It does not decay to our sector,
  requires no Coleman tunnelling, and supports stable matter structures
  with $\alpha'^{-1} = 139$.  The two sectors are related by twin prime
  symmetry of the Hecke eigenvalue analysis and correspond to two independent
  sets of modular forms (Set A and Set B).

  \textbf{Status: [MOTIVATED CONJECTURE --- two independent algebraic sectors]}
  \end{abstract}
\fi

\section{Mirror Sector: Vacuum Stability}
\label{sec:vacuum_stability}

\begin{statusbox}{Status: [MOTIVATED CONJECTURE --- two independent algebraic sectors]}
This analysis is motivated by the numerical observation that Hecke eigenvalues at
primes $p=137$ (Set A, our sector) and $p=139$ (Set B, mirror sector) are both
globally unique in the range $50$--$300$.  Set A and Set B correspond to modular
forms at different levels and weights, suggesting two \emph{algebraically independent}
sectors rather than an excitation of a single sector.  All conclusions beyond
the numerical observations are conjectural at this stage.
\end{statusbox}

\subsection{The Effective Potential}

The UBT effective potential for the vacuum mode number $n$ takes the one-loop form:
\begin{equation}
  V_{\rm eff}(n) = A\cdot n^2 - B\cdot n\ln(n),
  \label{eq:Veff}
\end{equation}
where:
\begin{itemize}
  \item $A$ is the quadratic coefficient from the kinetic term (positive; sets energy scale),
  \item $B$ is the logarithmic correction from one-loop running (positive; destabilises).
\end{itemize}

The minimum of $V_{\rm eff}(n)$ at $n = n^*$ satisfies $V_{\rm eff}'(n^*) = 0$:
\begin{equation}
  2A\cdot n^* = B\cdot (\ln n^* + 1)
  \;\Longrightarrow\;
  n^* = \frac{B}{2A}\left(\ln n^* + 1\right).
  \label{eq:minimum_condition}
\end{equation}
The observed value $n^* = 137$ gives the calibration relation
$B_0 = 2A\cdot 137 / (\ln 137 + 1)$.

\subsection{Numerical Analysis: $V_{\rm eff}(137)$ vs $V_{\rm eff}(139)$}

We evaluate $V_{\rm eff}$ at the twin prime values using the calibration
$B = B_{137} = 2A \cdot 137/(\ln 137 + 1)$.
With $A=1$ (setting energy units) and $\ln 137 \approx 4.920$:
\begin{equation}
  B_{137} = \frac{2 \times 137}{4.920 + 1} = \frac{274}{5.920} \approx 46.284.
  \label{eq:B137}
\end{equation}

\subsubsection{At $n = 137$:}
\begin{align}
  V_{\rm eff}(137) &= 137^2 - 46.284 \cdot 137 \cdot \ln(137) \notag\\
  &= 18769 - 46.284 \times 137 \times 4.920 \notag\\
  &= 18769 - 46.284 \times 674.04 \notag\\
  &= 18769 - 31197.1 \notag\\
  &\approx -12428.1.
  \label{eq:Veff137}
\end{align}

\subsubsection{At $n = 138$:}
\begin{align}
  V_{\rm eff}(138) &= 138^2 - 46.284 \cdot 138 \cdot \ln(138) \notag\\
  &= 19044 - 46.284 \times 138 \times 4.927 \notag\\
  &= 19044 - 46.284 \times 679.93 \notag\\
  &= 19044 - 31471.3 \notag\\
  &\approx -12427.3.
  \label{eq:Veff138}
\end{align}

\subsubsection{At $n = 139$:}
\begin{align}
  V_{\rm eff}(139) &= 139^2 - 46.284 \cdot 139 \cdot \ln(139) \notag\\
  &= 19321 - 46.284 \times 139 \times 4.934 \notag\\
  &= 19321 - 46.284 \times 685.83 \notag\\
  &= 19321 - 31745.8 \notag\\
  &\approx -12424.8.
  \label{eq:Veff139}
\end{align}

\begin{resultbox}
\textbf{[NUMERICAL OBSERVATION]:}
$V_{\rm eff}(137) \approx -12428.1 < V_{\rm eff}(139) \approx -12424.8$,
with $n^* = 137$ as the global minimum when $B = B_{137}$.

\medskip
\textbf{Important}: $n = 139$ is \emph{not} a local minimum of $V_{\rm eff}$
calibrated to our sector.  The discrete test confirms:
$V_{\rm eff}(138) \approx -12427.3 < V_{\rm eff}(139) \approx -12424.8$,
so $n=139$ lies on the ascending branch above the minimum.  The earlier
claim that $V''(139) > 0$ implies a local minimum was an error:
$V''(n) > 0$ holds for all $n > B/(2A) \approx 23$ since the potential
is convex in this region.
\end{resultbox}

\subsection{The Correct Interpretation: Two Independent Vacua}

The question is not \emph{``is $n=139$ a local minimum?''} but rather
\emph{``can $n^*=139$ be the minimum of a physically valid branch?''}.

The answer is yes.  The condition $V'_{\rm eff}(n^*) = 0$ gives:
\begin{equation}
  B_{n^*} = \frac{2A \cdot n^*}{\ln n^* + 1}.
\end{equation}
For the two twin primes:
\begin{align}
  B_{137} &= \frac{2 \times 137}{\ln 137 + 1} \approx 46.284, \\
  B_{139} &= \frac{2 \times 139}{\ln 139 + 1} \approx 46.845.
\end{align}
The relative difference is:
\begin{equation}
  \frac{\Delta B}{B} = \frac{B_{139} - B_{137}}{B_{137}} \approx 1.21\%.
\end{equation}

\begin{conjecturebox}
\textbf{[MOTIVATED CONJECTURE --- Two Independent Sectors]:}
The mirror sector is a branch of UBT with loop-correction coefficient
$B' = B_{139} \approx 46.845$, for which $n^* = 139$ is the \emph{global}
minimum.  This branch is \textbf{fully stable} --- it is not a false vacuum,
requires no tunnelling mechanism, and does not decay.

The two branches correspond to the two independent sets of Hecke modular forms:
\begin{itemize}
  \item Set A (levels 76, 7, 208): global minimum at $n^* = 137$, $B = B_{137}$
  \item Set B (levels 195, 50, 54): global minimum at $n^* = 139$, $B = B_{139}$
\end{itemize}
The algebraic independence of Set A and Set B (different levels, different weights,
different Hecke eigenvalues) is consistent with two autonomous sectors, not with
an excitation of a single sector.
\end{conjecturebox}

\subsection{Stability of the Mirror Sector}

Since $n^* = 139$ is the \emph{global} minimum of the mirror branch
(with $B = B_{139}$), no decay mechanism is required or expected.
Coleman tunnelling does not apply: there is no false vacuum.
The mirror sector is \textbf{permanently stable}.

The two sectors coexist as independent solutions of the UBT algebraic
structure.  Their interaction is purely gravitational (mirror matter couples
to our sector only via the $\phi$-projection, which is gravitational).

\subsection{Summary}

\begin{center}
\renewcommand{\arraystretch}{1.3}
\begin{tabular}{lll}
\toprule
\textbf{Result} & \textbf{Status} & \textbf{Notes} \\
\midrule
$V_{\rm eff}(137) < V_{\rm eff}(139)$ (same $B$) & \textbf{[NUMERICAL OBSERVATION]} &
  Calibrated to our sector \\
$n=139$ is not a local min of $V_{B_{137}}$ & \textbf{[PROVED NUMERICALLY]} &
  Discrete test: $V(138) < V(139)$ \\
$n^*=139$ is global min of branch $B_{139}$ & \textbf{[MOTIVATED CONJECTURE]} &
  $\Delta B/B = 1.21\%$ \\
Mirror sector is fully stable & \textbf{[FOLLOWS FROM CONJECTURE]} &
  No Coleman tunnelling needed \\
Set A and Set B are algebraically independent & \textbf{[NUMERICAL OBSERVATION]} &
  Different levels and weights \\
\bottomrule
\end{tabular}
\end{center}

\noindent\textbf{Date}: 2026-03-07 \\
\textbf{Author}: David Jaroš

\begin{thebibliography}{9}
\bibitem{MirrorSector}
R.\ Foot and R.\ R.\ Volkas, \textit{Neutrino physics and the mirror world:
How exact parity symmetry explains the solar neutrino deficit, the atmospheric
neutrino anomaly and the LSND experiment},
Phys.\ Rev.\ D \textbf{52}, 6595 (1995).
\end{thebibliography}

\ifdefined\INCLUDEMODE
\else
\end{document}
\fi
